Tato podsekce shrnuje praktické kroky a kontrolní body zaměřené na ochranu osobních a interních dat při nasazení AI služeb na úrovni fakulty/univerzity. Následující doporučení jsou navržena jako pracovní průvodce pro garanty předmětů, IT a právní/GDPR kontakty; konečná právní a procesní rozhodnutí je nutné konzultovat s interním právním oddělením a specialisty na ochranu osobních údajů (general background knowledge).

Poznámka k nástrojům ve výuce — ilustrativní příklady:
\begin{itemize}
  \item Některé edukační nástroje jsou součástí komerčních balíků (např. Minecraft Education v rámci Microsoft 365 A3/A5), což má vliv na smluvní podmínky a zpracování dat \cite{kb_ai_101_tipu_pdf_04e4a115_page_15_2}.
  \item Funkce založené na zpracování obrazu (OCR, extrakce textu z obrázků) mohou konvertovat nahrané skeny či screenshoty do zpracovatelného textu; to je relevantní při sdílení dokumentů obsahujících osobní údaje \cite{kb_ai_101_tipu_pdf_04e4a115_page_53_2}.
  \item Některé edukační AI funkce shromažďují a vyhodnocují výkon/student‑level metriky (např. nástroje generující příklady a analyzující postupy při matematice), což zvyšuje potřebu jasného smluvního ošetření a ochrany dat \cite{kb_ai_101_tipu_pdf_04e4a115_page_8_1}.
\end{itemize}

Co zkontrolovat v DPA (Data Processing Agreement) — praktické body (general background knowledge)
\begin{itemize}
  \item Role a odpovědnosti: jasné vymezení, kdo je správcem (controller) a kdo zpracovatelem (processor), včetně sub‑processorů a práv na audity.
  \item Účel zpracování: omezení na konkrétní účely výuky/administrace, zákaz sekundárního komerčního využití studentských dat bez souhlasu.
  \item Doba retence a pravidla mazání: termíny uchovávání, procesy pro úplné odstranění dat na vyžádání a mechanismy pro obnovu po chybě.
  \item Data residency a lokace zpracování: požadavek na zpracování/uložení v jurisdikci EU/EHP nebo jiná specifika podle interních pravidel fakulty.
  \item Šifrování a zabezpečení: povinnost používat šifrování v přenosu (TLS) a šifrování v klidu; volitelné zákaznické (customer‑managed) klíče jsou výhodou.
  \item Incidentní řízení: SLA pro hlášení úniku dat, povinnost notifikace do interních kontaktních bodů a lhůty pro oznámení.
  \item Podmínky sub‑processorů a auditu: seznam povolených sub‑processorů, právo na bezpečnostní audity a penále/odpovědnost za porušení.
\end{itemize}

Data residency — možnosti a implikace (general background knowledge)
\begin{itemize}
  \item Možné varianty: 1) úložiště a zpracování v EU/EHP, 2) regionální cloud s možnostmi smluvního omezení lokace, 3) globální cloud bez záruk residency. Volba ovlivní právní compliance a provozní latenci.
  \item Praktické pravidlo: pro citlivá studentská data preferovat EU‑based residency nebo on‑prem/hybridní řešení; u veřejně dostupných studijních materiálů lze zvážit širší varianty s odpovídajícími zárukami.
\end{itemize}

Šifrování a technická opatření (general background knowledge)
\begin{itemize}
  \item Požadujte šifrování v přenosu (TLS 1.2+/TLS 1.3) a šifrování v klidu (AES‑256 nebo ekvivalent).
  \item Zvažte customer‑managed keys (CMK) pro klíčovou kontrolu, rotaci klíčů a možnost odnětí přístupu v případě incidentu.
  \item U RAG scén a indexace interních zdrojů logujte retrieval události a volání modelu pro forenzní sledovatelnost (general background knowledge).
\end{itemize}

Praktický workflow pro sdílení citlivých materiálů v kurzu — krok po kroku (general background knowledge)
\begin{enumerate}
  \item Klasifikujte data používaná v kurzu (veřejné, interní, citlivé/osobní).
  \item Minimalizujte sběr: sbírejte pouze nezbytná pole (pseudonymizujte tam, kde lze).
  \item Proveďte DPIA (posouzení dopadů na ochranu osobních údajů) před zahájením rozsáhlejšího nasazení AI nástroje.
  \item Uzavřete nebo ověřte DPA se službou, zkontrolujte residency, šifrování a sub‑processor pravidla.
  \item Technicky nastavte přístup: SSO/IdP integrace, role‑based access, omezení exportů dat a auditní logování.
  \item Testujte PoC: ověřte mazání dat, export logů, chování při incidentu a chování extrakčních funkcí (OCR atp.).
  \item Do sylabu a komunikace s studenty jasně specifikujte pravidla sdílení dat a jak budou jejich výstupy zpracovávány (general background knowledge).
\end{enumerate}

Kontrolní checklist pro rychlý PoC (general background knowledge)
\begin{itemize}
  \item Existuje podepsaná DPA pokrývající všechny plánované datové toky?
  \item Je lokace uložených dat (data residency) v souladu s interní politikou?
  \item Jsou zapnutá šifrování v přenosu i v klidu a je dostupná možnost CMK?
  \item Jsou nastavena auditní logy pro přístupy k interním dokumentům a volání modelu?
  \item Proběhla DPIA a byla schválena příslušným orgánem/fakultním GDPR kontaktem?
\end{itemize}

Závěrečné doporučení pro garanty a IT (general background knowledge):
\begin{itemize}
  \item Zapojte právní/GDPR a IT/security týmy již ve fázi výběru řešení; smluvní zajištění a technické konfigurace ovlivní možnosti využití nástroje ve výuce.
  \item Piloty navrhujte s omezeným dosahem dat (pseudonymizace, anonymizace) a jasným plánem rozšíření až po ověření bezpečnosti a souladu.
  \item Dokumentujte všechna rozhodnutí o residency, DPA a retenčních politikách a zahrňte je do sylabů a komunikace ke studentům.
\end{itemize}